\section{讨论}
\label{sec:disc}

\subsection{收益出现的位置}

分层结果支持对该方法的一种特定解读。GAST 的相对优势并非均匀分布，而是随可用证据减少而扩大：相对 MoFlow，历史完整时 $\mathrm{minFDE}_{20}$ 降低 5.4\%，缺失三帧时降低 11.2\%，五帧时降低 14.8\%，六帧时降低 22.6\%；$\mathrm{minADE}_{20}$ 上对应为 6.7\% 至 25.9\%。当仅剩一个观测时该趋势反转，GAST 落后 MoFlow 约 4\%。按照证据条件化的设计意图，收益应集中在缺失较多的条件下，因为此时“哪些帧可信、最后一个状态有多旧”对预测的影响最大；观测完整时证据状态退化为恒等，模型与普通完整历史结构无异。观测到的形态与该预期相符。

误差随证据年龄的增长进一步说明问题来源。锚点间隔从 0 增至 7 时，GAST 的 $\mathrm{minFDE}_{20}$ 增至 4.0 倍；末端缺失组的误差比非末端组高 61.4\%。这两项都不是模型间比较，而是同一模型在不同证据条件下的行为，它们表明误差的主要来源是锚点状态的陈旧程度，而非场景本身。

但必须强调，以上只是与机制一致的证据，而非机制归因。本文未报告组件级消融，因此无法把收益分配给逐帧几何证据编码、可靠度调制的时间聚合与锚点条件化解码中的任何一个；同样无法排除模型容量、训练细节或归纳偏置带来的贡献。严格的归因需要在同一训练协议下逐组件开关的四臂实验，本文将其列为后续工作。

\subsection{场景异质性}

逐场景结果并不一致：GAST 在四个场景上取得最优 $\mathrm{minADE}_{20}$、在三个场景上取得最优 $\mathrm{minFDE}_{20}$，但在 HOTEL-M 上落后 AgentFormer（0.2541/0.3769 对 0.2083/0.3316），在 ZARA1-M 的 $\mathrm{minFDE}_{20}$ 上落后 MoFlow（0.4128 对 0.4057）。本文可以给出一个假设，尽管该基准无法对其进行干净的检验。HOTEL-M 与 ZARA1-M 的运动较慢、较规则，且受场景约束较强，从最后观测出发的单模态外推已经足够准确；以在多个不同目的地上分配概率质量为目标的目标条件化解码器，在此类场景中增加方差而不增加精度。ETH-M 与 UNIV-M 的运动更异质、方向变化更频繁，目的地本身更不确定，显式的终点假设因而有帮助。若该解读成立，解码器应依据场景级或样本级的终点不确定性估计进行条件化，而不是统一施加。本文认为这是缩小场景差距最有希望的方向，但当前证据只是相关性。

\subsection{覆盖度与排序}

best-of-20 覆盖度与 top-1 排序之间的落差是本研究中值得注意的现象。GAST 在 best-of-20 上相对 MoFlow 降低 9.4\%/7.6\%，但在概率最高的单一模式上两者基本相当（$\mathrm{ADE@1}$ 优 0.1\%，$\mathrm{FDE@1}$ 差 3.4\%）。换言之，收益主要来自候选集合对真实未来的覆盖能力提升，而非模式评分对最优候选的识别能力提升。

这一落差对使用方式有直接影响。best-of-20 适用于对候选集合做保守规划或集成的场合；若下游规划器只消费 top-1，本方法的优势将收窄至噪声量级。要同时改善排序，需要在训练目标中显式优化模式概率与轨迹质量的对应关系，例如引入排序损失或对模式评分做校准，而不是仅依赖轨迹级分类标签。本文未对此作进一步处理，并将其列为开放问题。

\subsection{局限性}

\textbf{单随机种子。}所有数值均来自每个场景的一次运行，没有置信区间。开发过程中同一配置在 UNIV-M 上的两次运行曾相差 4.4\%，因此单场景上几个百分点的差异（例如 ZARA1-M 上 1.8\% 的 $\mathrm{minFDE}_{20}$ 差距）处于噪声范围内。宏平均上 9.4\%/7.6\% 的领先大于该波动，但仍应视为单种子结论；在声称稳定改进之前，需要多种子研究与按轨迹聚类的配对自助重采样。

\textbf{协议依赖。}TrajGap-Bench 是基于合成掩码的受控压力测试，其缺失分布不对应任何真实跟踪器，且基准限定于五个 ETH/UCY 场景与 8 帧历史。真实感知输出与其他数据集（如 SDD 或 Argoverse）上的外部效度尚未验证。

\textbf{测试集依赖性。}测试划分中包含同一潜在轨迹在多个掩码下的多次出现。本文的聚合把样本视为独立，这会高估任何显著性陈述的有效样本量。发布的逐样本诊断包含源轨迹标识，正是为了便于后续进行聚类推断。

\textbf{缺少组件级消融。}本文报告的是完整模型与外部基线的对比，未逐一移除逐帧几何证据编码、可靠度调制的时间聚合与锚点条件化解码。因此方法内部的贡献分配尚不清楚，也无法判断某一模块是否可以被更简单的结构替代。

\textbf{缺失的对比。}本文未在直接预测表中报告两阶段插补基线，因为该设置向模型暴露更多信息、必须单独报告；因此这是对比的缺失，而非阴性结果。本文亦未报告推理时延，而这对部署是重要的。

\subsection{对基准设计的启示}

协议中有两项设计被证明是承重的。将 Easy 与 Hard 合并为单一 Mixed 划分，使模型无法在“自己恰好最强的难度”上被选中，并使缺失计数分层成为诊断而非榜单。把检查点固定为最终预测头的验证集 $\mathrm{minFDE}_{20}$，消除了一个原本会因论文而异的自由度。两项选择都易于采纳，且据本文经验会改变结论：在按难度分别选点的口径下，不同检查点之间的排序可以发生改变，因而跨论文比较必须依赖统一的选点规则。本文建议后续缺失历史预测工作报告至少缺失帧数与末端缺失两项分层，因为聚合值可能掩盖文中七帧处记录的反转。
